\chapter{Convergence results for Ricci flow}
\label{chap:convergence-results}


This chapter records geometric, Gromov--Hausdorff, and blow-up
convergence results for Riemannian manifolds and Ricci flows, including
the splitting result at infinity.



\section{Geometric convergence of Riemannian manifolds}

\begin{definition}[$\delta$-regular point]\label{def:delta-regular-point}
\uses{def:riemannian-metric}
\notready
\label{deltareg}

Let $(U,g)$ be a Riemannian manifold. Let $\delta>0$ be given. We
say that $p\in U$ is a {\em $\delta$-regular}
point if for every $r'<\delta$ the
metric ball $B(p,r')$ has compact closure in $U$. Equivalently, $p$
is $\delta$-regular if the exponential mapping at $p$ is defined on
the open ball of radius $\delta$ centered at the origin in $T_pU$,
i.e., if each geodesic ray emanating from $p$ extends to a geodesic
defined on $[0,\delta)$. We denote by $\operatorname{Reg}_\delta(U,g)$ the
subset of $\delta$-regular points in $(U,g)$. For any $x\in \operatorname{Reg}_\delta(U,g)$ we denote by $\operatorname{Reg}_\delta(U,g,x)$ the
connected component of $\operatorname{Reg}_\delta(U,g)$ containing $x$.
\end{definition}


\begin{lemma}[Regular points form a closed set]\label{lem:reg-delta-closed}
\uses{def:delta-regular-point}
\notready

 $\operatorname{Reg}_\delta(U,g)$  is a closed subset of $U$.
\end{lemma}

\begin{proof}
\sketch
Suppose that $p_n$ converges to $p$ as $n$ tends to $\infty$ and
suppose that $p_n\in \operatorname{Reg}_\delta(U,g)$ for all $n$. Fix
$r'<\delta$ and consider the
 ball $B(p,r')$. For all $n$ sufficiently large, this
ball is contained in $B(p_n,(\delta+r')/2)$, and hence has compact
closure.
\end{proof}




\begin{definition}[Geometric limit of Riemannian manifolds]\label{def:geometric-limit}
\uses{def:delta-regular-point, def:riemannian-metric}
\notready
\label{smoothconv}

For each $k$
let $(U_k,g_k,x_k)$ be a based, connected Riemannian manifold. A
{\sl geometric limit} of the sequence $\{U_k,g_k,x_k\}_{k=0}^\infty$
is a based, connected Riemannian manifold
$(U_\infty,g_\infty,x_\infty)$
 with the extra data:
\begin{enumerate}
\item[(1)] An increasing sequence $V_k\subset U_\infty$ of connected open subsets
of $U_\infty$ whose union is $U_\infty$ and which satisfy the
following for all $k$:
\begin{enumerate}
\item the closure $\overline V_k$ is compact, \item
$\overline V_k\subset V_{k+1}$,  \item $V_k$ contains $x_\infty$.
\end{enumerate}
\item For each $k\ge 0$ a smooth embedding $\varphi_k\colon (V_k,x_\infty)\to
(U_k,x_k)$ with the properties that: \begin{enumerate} \item
$\lim_{k\rightarrow \infty}\varphi_k^*g_k=g_\infty$, where the limit
is in the uniform $C^\infty$-topology on compact subsets of
$U_\infty$.
\item For any $\delta>0$ and any $R<\infty$ for all $k$ sufficiently large,
$x_k\in \operatorname{Reg}_\delta(U_k,g_k)$ and for any $\ell\ge k$ the image
$\varphi_\ell(V_k)$ contains $B(x_\ell,R)\cap\operatorname{Reg}_\delta(U_\ell,g_\ell,x_\ell)$.
\end{enumerate}
\end{enumerate}

We also say that the sequence {\em converges geometrically
to} $(U_\infty,g_\infty,x_\infty)$ if
there exist $(V_k,\varphi_k)$ as required in the above definition.
We also say that $(U_\infty,g_\infty,x_\infty)$ is the {\em
geometric limit} of the sequence.

More generally, given $(U_\infty,g_\infty,x_\infty)$, a sequence of open
subsets and $\{V_k\}_{k=1}^\infty$ satisfying (1) above, and smooth maps
$\varphi_k\colon V_k\to U_k$ satisfying (2a) above, we say that
$(U_\infty,g_\infty,x_\infty)$ is a {\em partial geometric
limit} of the sequence.
\end{definition}


\begin{remark}[Geometric limits and incompleteness]\label{rem:geometric-limit-incomplete-condition}
\uses{def:geometric-limit}
\notready

Conditions (1) and (2a) in the definition above also appear in the definition
in the case of complete limits. It is Condition (2b) that is extra in this
incomplete case. It says that once $k$ is sufficiently large then the image
$\varphi_\ell(V_k)$ contains all points satisfying two conditions: they are at
most a given bounded distance from $x_\ell$, and also they are at least a fixed
distance from the boundary of $U_\ell$.
\end{remark}




\begin{lemma}[Uniqueness of geometric limits]\label{lem:geometric-limit-uniqueness}
\uses{def:geometric-limit}
\notready

The geometric limit of a sequence $(U_k,g_k,x_k)$ is unique up to
based isometry.
\end{lemma}

\begin{proof}
\sketch
\uses{def:geometric-limit,def:delta-regular-point}
Suppose that we have two geometric limits
$(U_\infty,g_\infty,x_\infty)$ and
$(U'_\infty,g'_\infty,x'_\infty)$. Let $\{V_k,\varphi_k\}$ and
$\{V'_k,\varphi_k'\}$ be the sequences of open subsets and maps as
required by the definition of the limit.

Fix $k$. Since $V_k$ is connected and has compact closure, there are
$R<\infty$ and $\delta>0$ such that $V_k\subset B(x_\infty,R)\cap
\operatorname{Reg}_\delta(U_\infty,g_\infty,x_\infty)$. Let $x$ be contained
in the closure of $V_k$. Then by the triangle inequality the closed
ball $\overline{B(x,\delta/3)}$ is contained in
$B(x_\infty,R+\delta)\cap\operatorname{Reg}_{\delta/2}(U_\infty,g_\infty,x_\infty)$. Since the union of
these closed balls as $x$ ranges over $\overline V_k$ is a compact
set, for all $\ell$ sufficiently large, the restriction of
$\varphi_\ell^*g_\ell$ to the union of these balls is close to the
restriction of $g_\infty$ to the same subset. In particular, for all
$\ell$ sufficiently large and any $x\in V_k$ we see that
$\varphi_\ell\left(B(x,\delta/3)\right)$ contains
$B(\varphi_\ell(x),\delta/4)$. Thus, for all $\ell$ sufficiently
large $\varphi_\ell(V_k)\subset B(x_\ell,R+2\delta)\cap \operatorname{Reg}_{\delta/4}(U_\ell,g_\ell,x_\ell)$. This implies that, for given
$k$, for all $\ell$ sufficiently large $\varphi_\ell(V_k)\subset
\varphi'_\ell(V'_\ell)$. Of course, $(\varphi'_\ell)^{-1}\circ
\varphi_\ell(x_\infty)=x_\infty'$. Fix $k$ and pass to a subsequence
of $\ell$, such that as $\ell\rightarrow\infty$, the compositions
$(\varphi'_\ell)^{-1}\circ\left(\varphi_\ell|_{V_k}\right)\colon
V_k\to U'_\infty$ converge to a base-point preserving isometric
embedding of $V_k$ into $U'_\infty$. Clearly, as we pass from $k$ to
$k'>k$ and take a further subsequence of $\ell$ these limiting
isometric embeddings are compatible. Their union is then a
base-point preserving isometric embedding of $U_\infty$ into
$U'_\infty$.

The last thing we need to see is that the embedding of $U_\infty$ into
$U'_\infty$ constructed in the previous paragraph is onto. For each $n$ we have
$\overline V'_n\subset V'_{n+1}$. Since $\overline V'_n$ is compact and
connected, it follows that there are $R<\infty$ and $\delta>0$ (depending on
$n$) such that $\overline V'_n\subset B(x'_\infty,R)\cap \operatorname{Reg}_\delta(V_{n+1},g'_\infty,x_\infty')$. Since $V'_{n+1}$ has compact closure
in $U'_\infty$, as $\ell$ tends to $\infty$ the metrics
$(\varphi'_\ell)^*g_\ell$ converge uniformly on $V_{n+1}$ to
$g'_\infty|_{V_{n+1}}$. This means that there are $R'<\infty$ and $\delta'>0$
(depending on $n)$ such that for all $\ell$ sufficiently large,
$\varphi'_\ell(V_n)\subset B(x_k,R')\cap \operatorname{Reg}_{\delta'}(U_\ell,g_\ell,x_\ell)$. This implies that for all $k$
sufficiently large and any $\ell\ge k$ the image $\varphi'_\ell(V'_n)$ is
contained in the image of $\varphi_\ell(V_k)$. Hence, for all $k$  sufficiently
large and any $\ell\ge k$ we have
$V'_n\subset(\varphi'_\ell)^{-1}(\varphi_\ell(V_k))$. Hence, the isometric
embedding $U_\infty\to U'_\infty$ constructed above contains $V'_n$. Since this
is true for every $n$, it follows that this isometric embedding is in fact an
isometry $U_\infty\to U'_\infty$.
\end{proof}








\begin{theorem}[Existence of geometric limits]\label{thm:geometric-limit-existence}
\uses{def:delta-regular-point}
\notready
\label{basicconv}

Suppose that $\{(U_k,g_k,x_k)\}_{k=1}^\infty$ is a sequence of
based, connected, $n$-dimensional Riemannian manifolds. In addition,
suppose the following:

\begin{enumerate}
\item[(1)] There is $\delta>0$ such that $x_k\in \operatorname{Reg}_\delta(U_k,g_k)$ for all $k$.
\item[(2)] For each $R<\infty$ and $\delta>0$ there is a constant $V(R,\delta)<\infty$ such
that $\Vol(B(x_k,R)\cap \operatorname{Reg}_\delta(U_k,x_k))\le
V(R,\delta)$ for all $k$ sufficiently large.
\item[(3)] For each non-negative integer $\ell$, each $\delta>0$, and each $R<\infty$,
 there is a constant
$C(\ell,\delta,R)$ such that for every $k$ sufficiently large we
have
$$|\nabla^\ell \Rm(g_k)|\le C(\ell,\delta,R)$$
on all of $B(x_k,R)\cap\operatorname{Reg}_\delta(U_k,g_k)$.
\item[(4)] For every $R<\infty$ there are $r_0>0$ and $\kappa>0$ such that
for every $k$ sufficiently large, for every $\delta\le r_0$ and
every $x\in B(x_k,R)\cap \operatorname{Reg}_\delta(U_k,g_k,x_k)$ the volume
of the metric ball $B(x,\delta)\subset U_k$ is at least
$\kappa\delta^n$.
\end{enumerate}
Then, after passing to a subsequence, there exists a based
Riemannian manifold $(U_\infty,g_\infty,x_\infty)$ that is a
geometric limit of the sequence $\{(U_k,g_k,x_k)\}_{k=1}^\infty$.
\end{theorem}



\begin{proof}
\sketch
\uses{thm:volume-injectivity-radius,def:delta-regular-point,def:geometric-limit,lem:gaussian-coordinates-convergence}
Fix $R<\infty$ and $\delta>0$. Let
$$X(\delta,R)=B(x_k,R)\cap
\operatorname{Reg}_{2\delta}(U_k,g_k,x_k).$$
 From the non-collapsing assumption
and the curvature bound assumption if follows from
Theorem~\ref{thm:volume-injectivity-radius} that there is a uniform positive lower bound
(independent of $k$) to the injectivity radius of every point in
$X(\delta,R)$. Fix $0<\delta'\le \min(r_0,\delta/2)$ much less
than this injectivity radius. We also choose $\delta'>0$
sufficiently small so that any ball of radius $2\delta'$ in
$B(x_k,R+\delta)\cap \operatorname{Reg}_\delta(U_k,g_k,x_k)$ is geodesically
convex. (This is possible because of the curvature bound.)  We cover
$X(\delta,R)$ by balls $B_1',\ldots,B_N'$
 of radii $\delta'/2$ centered at points of $X(\delta,R)$ with the
property that the sub-balls of radius $\delta'/4$ are disjoint. We
denote by $B_i'\subset B_i\subset \widetilde B_i$ the metric balls
with the same center and radii $\delta'/2$, $\delta'$, and
$2\delta'$ respectively.
 Notice that each of the balls $\widetilde B_i$  is contained in $B(x_k,R+\delta)\cap \operatorname{Reg}_{\delta}(U_k,g_k,x_k)$. Because $\delta'\le r_0$, because $\Vol\,B(x_k,R+\delta)$ is bounded independent of $k$, and because the
concentric balls of radius $\delta'/4$ are disjoint, there is a uniform bound
(independent of $k$) to the number of such balls. Passing to a subsequence we
can assume that the number of balls in these coverings is the same for all $k$.
We number them $\widetilde B_1,\ldots, \widetilde B_N$. Next, using the
exponential mapping at the central point, identify each of these balls with the
ball of radius $2\delta'$ in $\Ar ^n$. By passing to a further subsequence we
can arrange that the metrics on each $\widetilde B_i$ converge uniformly. (This
uses the fact that the concentric balls of radius $2\delta\ge 4\delta'$ are
embedded in the $U_k$ by the exponential mapping.) Now we pass to a further
subsequence so that the distance between the centers of the balls converges,
and so that for any pair $\widetilde B_i$ and $\widetilde B_j$ for which the
limiting distance between their centers is less than $4\delta'$, the overlap
functions in the $U_k$ also converge. The limits of the overlap functions
defines a limiting equivalence relation on $\coprod_i\widetilde B_i$.

This allows us to form a limit manifold $\widehat U_\infty$. It is
the quotient of the disjoint union of the $\widetilde B_i$ with the
limit metrics under the limit equivalence relation. We set
$(U_\infty(\delta,R),g_\infty(\delta,R),x_\infty(\delta,R))$ equal
to the submanifold of $\widehat U_\infty$ that is the union of the
sub-balls $B_i\subset \widetilde B_i$ of radii $\delta'$. A standard
argument using partitions of unity and the geodesic convexity of the
balls $\widetilde B_i$ shows that, for all $k$ sufficiently large,
there are smooth embeddings $\varphi_k(\delta,R)\colon
U_\infty(\delta,R)\to B(x_k,R+\delta)\cap \operatorname{Reg}_\delta(U_k,g_k,x_k)$ sending $x_\infty(\delta,R)$ to $x_k$ and
converging as $k\rightarrow \infty$, uniformly in the
$C^\infty$-topology on each $B_i$, to the identity. Furthermore, the
images of each of these maps contains $B(x_k,R)\cap \operatorname{Reg}_{2\delta}(U_k,g_k,x_k)$; compare \cite{Cheeger}. Also, the pull
backs under these embeddings of the metrics $g_k$ converge uniformly
to $g_\infty(\delta,R)$.

Repeat the process with $R$ replaced by $2R$ and $\delta=\delta_1$
replaced by $\delta_2\le \delta_1/2$. This produces
$$\left(U_\infty(\delta_2,2R),g_\infty(\delta_2,2R),x_\infty(\delta_2,2R)\right)$$
and, for all $k$ sufficiently large, embeddings
$\varphi_k(\delta_2,2R)$ of this manifold into
$$B(x_k,2R+\delta_2)\cap \operatorname{Reg}_{\delta_2}(U_k,g_k,x_k).$$ Hence,
the image of these embeddings contains the images of the original
embeddings. The compositions $(\varphi_k(\delta_2,2R))^{-1}\circ
\varphi_k(\delta,R)$ converge to an isometric embedding
$$\left(U_\infty(\delta,R),g_\infty(\delta,R),x_\infty(\delta,R)\right)\to
\left(U_\infty(\delta_2,2R),g_\infty(\delta_2,2R),x_\infty(\delta_2,2R)\right).$$
Repeating  this construction infinitely often produces a manifold
$(U_\infty,g_\infty,x_\infty)$ which is written as an increasing union of open
subsets $V_k=U_\infty(\delta_k,2^kR)$, where the $\delta_k$ tend to zero as $k$
tends to $\infty$. For each $k$ the open subset $V_k$ has compact closure
contained in $V_{k+1}$. By taking a subsequence of the original sequence we
have maps $\varphi_k\colon V_k\to U_k$ so that (2a) in the definition of
geometric limits holds. Condition (2b) clearly holds by construction.
\end{proof}

\begin{lemma}[Convergence of balls in Gaussian coordinates]\label{lem:gaussian-coordinates-convergence}
\notready
Suppose that we have a sequence of $n$-dimensional balls $(B_k,h_k)$
of radius $r$ in Riemannian $n$-manifolds. Suppose that for each
$\ell$ there is a constant $C(\ell)$ such that for every $k$, we
have $|\nabla ^\ell \Rm(h_k)|\le C(\ell)$ throughout $B_k$.
Suppose also that for each $n$ the exponential mapping from the
tangent space at the center of $B_k$ induces a diffeomorphism from a
ball in the tangent space onto $B_k$. Then choosing an isometric
identification of the tangent spaces at the central points of the
$B_k$ with $\Ar^n$ and pulling back the metrics $h_k$ via the
exponential mapping to metrics $\tilde h_k$ on the ball $B$ of
radius $r$ in $\Ar^n$ gives us a family of metrics on $B$ that,
after passing to a subsequence, converge in the $C^\infty$-topology,
uniformly on compact subsets of $B$, to a limit.
\end{lemma}

\begin{lemma}[Complete partial limits are geometric limits]\label{lem:partial-limit-complete-is-limit}
\uses{def:geometric-limit}
\notready

Suppose that $(U_k,g_k,x_k)$ is a sequence of based Riemannian
manifolds and that there is a partial geometric limit
$(U_\infty,g_\infty,x_\infty)$ that is a complete Riemannian
manifold. Then this partial geometric limit is a geometric limit.
\end{lemma}




\begin{proof}
\sketch
\uses{def:geometric-limit,def:delta-regular-point}
Since the balls $B(x_\infty,R)$ have compact closure in $U_\infty$
and since
$$\operatorname{Reg}_\delta(U_\infty,g_\infty,x_\infty)=U_\infty$$ for
every $\delta>0$, it is easy to see that the extra condition, (2b), in
Definition~\ref{def:geometric-limit} is automatic in this case.
\end{proof}

\begin{theorem}[Geometric convergence for complete manifolds]\label{thm:manifold-convergence-complete-limit}
\uses{thm:geometric-limit-existence}
\notready
\label{mfdconv}

Let $\{(M_k,g_k,x_k)\}_{k=1}^\infty$ be a sequence  of connected,
based Riemannian manifolds. Suppose that:
\begin{enumerate}
\item[(1)] For every $A<\infty$ the ball
$B(x_k,A)$ has compact closure in $M_k$ for all $k$ sufficiently
large.
\item[(2)] For each integer $\ell\ge 0$ and each $A<\infty$ there is a constant
$C=C(\ell,A)$ such that for each $y_k\in B(x_k,A)$ we have
$$\left|\nabla^\ell\Rm(g_k)(y_k)\right|\le C$$
for all $k$ sufficiently large.
\item[(3)] Suppose also that there is a constant $\delta>0$ such that $\inj_{(M_k,g_k)}(x_k)\ge \delta$ for all $k$ sufficiently large.
\end{enumerate}
 Then after passing to a subsequence there is a geometric limit which
 is a complete Riemannian manifold.
\end{theorem}

\begin{proof}
\sketch
\uses{thm:bishop-gromov,thm:geometric-limit-existence}
By the curvature bounds, it follows from the Bishop-Gromov theorem
(Theorem~\ref{thm:bishop-gromov}) that for each $A<\infty$ there is a uniform bound
to the volumes of the balls $B(x_k,A)$ for all $k$ sufficiently large. It also
follows from the same result that  the uniform lower bound on the injectivity
radius at the central point implies that for each $A<\infty$ there is a uniform
lower bound for the injectivity radius on the entire ball $B(x_k,A)$, again for
$k$ sufficiently large. Given these two facts, it follows immediately from
Theorem~\ref{thm:geometric-limit-existence} that there is a geometric limit.

Since, for every $A<\infty$, the $B(x_k,A)$ have compact closure in
$M_k$ for all $k$ sufficiently large, it follows that for every
$A<\infty$ the ball $B(x_\infty,A)$ has compact closure in
$M_\infty$. This means that $(M_\infty,g_\infty)$ is complete.
\end{proof}

\begin{corollary}[Geometric convergence under non-collapsing]\label{cor:manifold-convergence-volume-noncollapsed}
\uses{thm:manifold-convergence-complete-limit}
\notready
\label{2ndmfdconv}

Suppose that $\{(M_k,g_k,x_k)\}_{k=1}^\infty$ is a sequence of based, connected
Riemannian manifolds. Suppose that the first two conditions in
Theorem~\ref{thm:manifold-convergence-complete-limit} hold and suppose also that there are constants $\kappa>0$
and $\delta>0$ such that $\Vol_{g_k}B(x_k,\delta)\ge \kappa\delta^n$ for
all $k$. Then after passing to a subsequence there is a geometric limit which
is a complete Riemannian manifold.
\end{corollary}

\begin{proof}
\sketch
\uses{thm:bishop-gromov,thm:volume-injectivity-radius,thm:manifold-convergence-complete-limit}
Let $A=\max(\delta^{-2},C(0,\delta))$, where $C(0,\delta)$ is the constant
given in the second condition in Theorem~\ref{thm:manifold-convergence-complete-limit}. Rescale, replacing the
Riemannian metric $g_k$ by $Ag_k$. Of course, the first condition of
Theorem~\ref{thm:manifold-convergence-complete-limit} still holds as does the second with different constants,
and we have $\left|\Rm_{Ag_k}(y_k)\right|\le 1$ for all $y_k\in
B_{Ag_k}(x_k,\sqrt{A}\delta)$. Also, $\Vol\,B_{Ag_k}(x_k,\sqrt{A}\delta)\ge \kappa (\sqrt{A}\delta)^n$. Thus, by the
Bishop-Gromov inequality (Theorem~\ref{thm:bishop-gromov}), we have $\Vol_{Ag_k}B(x_k,1)\ge \kappa/\Omega$ where
$$\Omega=\frac{V(\sqrt{A}\delta)}{(\sqrt{A}\delta)^nV(1)},$$
where $V(a)$ is the volume of the ball of radius $a$ in hyperbolic
$n$-space (the simply connected $n$-manifold of constant curvature
$-1$).
 Since
$\sqrt{A}\delta\ge 1$, this proves that for the rescaled manifolds the absolute
values of the sectional curvatures on $B_{Ag_k}(x_k,1)$ are bounded by $1$ and
the $\Vol_{Ag_k}B(x_k,1)$ are bounded below by a positive constant
independent of $k$.
 According to Theorem~\ref{thm:volume-injectivity-radius} the lower bound on the volume of the ball
of radius $1$ and the curvature bound on the ball of radius $1$ yield a uniform
positive lower bound $r>0$ for the injectivity radius of the rescaled manifolds
at $x_k$. Hence, the injectivity radii at the base points of the original
sequence are bounded below by $\delta/\sqrt{A}$. This means that the original
sequence of manifolds satisfies the third condition in Theorem~\ref{thm:manifold-convergence-complete-limit}.
Invoking this theorem gives the result.
\end{proof}


\subsection{Geometric convergence of manifolds in the case of Ricci
flow}

As the next theorem shows, because of Shi's theorem, it is much
easier to establish the geometric convergence manifolds in the
context of Ricci flows than in general.




\begin{theorem}[Geometric convergence of Ricci flow time-slices]\label{thm:ricci-flow-slice-convergence}
\uses{def:generalized-ricci-flow, def:time-slices-horizontal, def:compatible-embedding}
\notready
\label{sliceconv}

Suppose that $({\mathcal M}_k,G_k,x_k)$ is a sequence of based
generalized $n$-dimensional Ricci flows with ${\bf t}(x_k)=0$. Let
$(M_k,g_k)$ be the $0$ time-slice of $({\mathcal M}_k,G_k)$. Suppose
that for each $A<\infty$ there are constants $C(A)<\infty$ and
$\delta(A)>0$ such that for all $k$ sufficiently large the following
hold:
\begin{enumerate}
\item[(1)] the ball $B(x_k,0,A)$ has compact closure in $M_k$,
\item[(2)] there is an embedding $B(x_k,0,A)\times (-\delta(A),0]\to {\mathcal M}_k$ compatible with
the time function and with the vector field, \item[(3)]   $|\Rm|\le C(A)$
on the image of the embedding in the Item (2), and \item[(4)] there is $r_0>0$
and $\kappa>0$ such that $\Vol\,B(x_k,0,r_0)\ge \kappa r_0^n$ for all $k$
sufficiently large.
\end{enumerate} Then, after passing to a
subsequence, there is a geometric limit
$(M_\infty,g_\infty,x_\infty)$ of the $0$ time-slices
$(M_k,g_k,x_k)$. This limit is a complete Riemannian manifold.
\end{theorem}


\begin{proof}
\sketch
\uses{thm:manifold-convergence-complete-limit,thm:shi-derivative-estimates,cor:manifold-convergence-volume-noncollapsed}
The first condition in Theorem~\ref{thm:manifold-convergence-complete-limit} holds by our first
assumption. It is immediate from Shi's theorem (Theorem~\ref{thm:shi-derivative-estimates})
that the second condition of Theorem~\ref{thm:manifold-convergence-complete-limit} holds. The result
is then immediate from Corollary~\ref{cor:manifold-convergence-volume-noncollapsed}.
\end{proof}


\section{Geometric convergence of Ricci flows}




\begin{definition}[Partial geometric limit Ricci flow]\label{def:partial-geometric-limit-ricci-flow}
\uses{def:geometric-limit, def:generalized-ricci-flow, def:time-slices-horizontal, def:compatible-embedding}
\notready

Let $\{({\mathcal M}_k,G_k,x_k)\}_{k=1}^\infty$ be a sequence of based
generalized Ricci flows. We suppose that ${\bf t}(x_k)=0$ for all $k$ and we
denote by $(M_k,g_k)$ the time-slice of $({\mathcal M}_k,G_k)$. For some
$0<T\le \infty$, we say that a based Ricci flow
$(M_\infty,g_\infty(t),(x_\infty,0))$ defined for $t\in (-T,0]$ is a {\sl
partial geometric limit Ricci flow} if:
\begin{enumerate}
\item
There are open subsets $x_\infty\in V_1\subset V_2\subset \cdots \subset
M_\infty$ satisfying (1) of Definition~\ref{def:geometric-limit} with $M_\infty$ in
place of $U_\infty$,
\item there is a sequence $0<t_1<t_2<\cdots$ with  $\lim_{k\rightarrow\infty}t_k=T$, \item and maps
$$\widetilde\varphi_k\colon (V_k\times [-t_k,0])\to {\mathcal M}_k$$
compatible with time and the vector field
\end{enumerate}
such that the sequence of horizontal families of metrics
$\widetilde\varphi_k^*G_k$ converges uniformly on compact subsets of
$M_\infty\times (-T,0]$ in the $C^\infty$-topology to the horizontal
family of metrics $g_\infty(t)$ on $M_\infty\times (-T,0]$.
\end{definition}


\begin{definition}[Geometric limit Ricci flow]\label{def:geometric-limit-ricci-flow}
\uses{def:partial-geometric-limit-ricci-flow, def:geometric-limit}
\notready

For $0<T\le \infty$, if $(M_\infty,g_\infty(t),x_\infty),\ -T<t\le 0$, is a
partial geometric limit Ricci flow  of the based generalized Ricci flows
$({\mathcal M}_k,G_k,x_k)$ and if $(M_\infty,g_\infty(0),x_\infty)$ is a
geometric limit of the $0$ time-slices, then we say that the partial geometric
limit is a {\em geometric limit  Ricci flow defined on the time interval
$(-T,0]$}.
\end{definition}

\begin{proposition}[Existence of partial geometric limit flows]\label{prop:partial-flow-limit-existence}
\uses{def:geometric-limit, def:compatible-embedding, def:generalized-ricci-flow}
\notready
\label{partialflowlimit}

Hamilton's compactness construction is summarized in \cite{Hamiltonlimits}.
Fix constants $-\infty\le T'\le 0\le T\le \infty$ and suppose that $T'<T$. Let
$\{({\mathcal M}_k,G_k,x_k)\}_{k=1}^\infty$ be a sequence of based generalized
Ricci flows. Suppose that ${\bf t}(x_k)=0$ for all $k$, and denote by
$(M_k,g_k)$ the $0$ time-slice of $({\mathcal M}_k,G_k)$. Suppose that there is
a partial geometric limit $(M_\infty,g_\infty,x_\infty)$ for the
$(M_k,g_k,x_k)$ with open subsets $\{V_k\subset M_\infty\}$ and maps
$\varphi_k\colon V_k\to M_k$ as in Definition~\ref{def:geometric-limit}. Suppose that
for every compact subset $K\subset M_\infty$ and every compact interval
$I\subset (T',T)$ containing $0$, for all $k$ sufficiently large, there is an
embedding $\widetilde \varphi_k(K,I)\colon K\times I\to {\mathcal M}_k$
compatible with time and the vector field and extending the map $\varphi_k$ on
the $0$ time-slice. Suppose in addition that for every $k$ sufficiently large
there is a uniform bound (independent of $k$) to the norm of Riemann curvature
on the image of $\widetilde\varphi_k(K,I)$. Then after passing to a subsequence
the flows $\widetilde\varphi_k^*G_k$ converge  to a partial geometric limit
Ricci flow $g_\infty(t)$ defined for $t\in (T',T)$.
\end{proposition}

\begin{proof}
\sketch
\uses{thm:shi-derivative-estimates,def:compatible-embedding}
 Suppose that we have a partial geometric limit of the
time-zero slices as stated in the proposition. Fix a compact subset $K\subset
M_\infty$ and a compact sub-interval $I\subset (T',T)$. For all $k$
sufficiently large we have embeddings $\widetilde \varphi_k(K,I)$ as stated. We
consider the flows $g_k(K,I)(t)$ on $K\times I$ defined by pulling back the
horizontal metrics $G_k$ under the maps $\widetilde \varphi_k(K,I)$. These of
course satisfy the Ricci flow equation on $K\times I$. Furthermore, by
assumption the flows $g_k(K,I)(t)$ have uniformly bounded curvature. Then under
these hypothesis, Shi's theorem can be used to show that the curvatures of the
$g_k(K,I)$ are uniformly bounded $C^\infty$-topology. The basic computation
 done by Hamilton in \cite{Hamiltonlimits} shows that after passing
to a further subsequence, the Ricci flows $g_k(K,I)$ converge
uniformly in the $C^\infty$-topology to a limit flow on $K\times I$.
A standard diagonalization argument allows us to pass to a further
subsequence so that the pullbacks $\widetilde\varphi_k^*G_k$
converge uniformly in the $C^\infty$-topology on every compact
subset of $M_\infty\times (T',T)$. Of course, the limit satisfies
the Ricci flow equation.
\end{proof}




This `local' result leads immediately to the following result for
complete limits.

\begin{theorem}[Existence of geometric limit Ricci flows]\label{thm:ricci-flow-geometric-limit-existence}
\uses{thm:ricci-flow-slice-convergence, prop:partial-flow-limit-existence, def:compatible-embedding}
\notready
\label{flowlimit}

Fix $-\infty\le T'\le 0\le T\le \infty$ with $T'<T$. Let $\{({\mathcal
M}_k,G_k,x_k)\}_{k=1}^\infty$ be a sequence of based generalized Ricci flows.
Suppose that ${\bf t}(x_k)=0$ for all $k$, and denote by $(M_k,g_k)$ the $0$
time-slice of $({\mathcal M}_k,G_k)$. Suppose that for each $A<\infty$ and each
compact interval $I\subset (T',T)$ containing $0$ there is a constant $C(A,I)$
such that the following hold for all $k$ sufficiently large:
\begin{enumerate}
\item[(1)]  the ball
$B_{g_k}(x_k,0,A)$ has compact closure in $M_k$, \item[(2)]  there is an
embedding $B_{g_k}(x_k,0,A)\times I\to {\mathcal M}_k$ compatible with time and
with the vector field,
\item[(3)] the norms of the Riemann curvature of $G_k$ on the image of the embedding in the
previous item are  bounded by $C(A,I)$, and
\item[(4)] there is $r_0>0$ and $\kappa>0$ with $\Vol\,B(x_k,0,r_0)\ge \kappa r_0^n$
for all $k$ sufficiently large.
\end{enumerate}
 Then after passing to a subsequence there is a
flow  $(M_\infty,g_\infty(t),(x_\infty,0))$ which is the geometric
limit. It is a solution to the Ricci flow equation defined for $t\in
(T',T)$. For every $t\in (T',T)$ the Riemannian manifold
$(M_\infty,g_\infty(t))$ is complete.
\end{theorem}


\begin{proof}
\sketch
\uses{thm:ricci-flow-slice-convergence,prop:partial-flow-limit-existence}
By Theorem~\ref{thm:ricci-flow-slice-convergence} there is a geometric limit $(M_\infty,g_\infty(0))$
of the $0$ time-slices, and the limit  is a complete Riemannian manifold. Then
by Proposition~\ref{prop:partial-flow-limit-existence} there is a geometric limit flow defined
on the time interval $(T',T)$. Since for every $t\in (T',T)$ there is a compact
interval $I$ containing $0$ and $t$, it follows that the Riemann curvature of
the limit is bounded on $M_\infty\times I$. This means that the metrics
$g_\infty(0)$ and $g_\infty(t)$ are commensurable with each other. Since
$g_\infty(0)$ is complete so is $g_\infty(t)$.
\end{proof}


\begin{corollary}[Extension of Ricci flow with bounded curvature]\label{cor:ricci-flow-extension-bounded-curvature}
\uses{def:ricci-flow}
\notready

Suppose that $(U,g(t)),\ 0\le t<T<\infty$,\ is a Ricci flow. Suppose
that $|\Rm(x,t)|$ is bounded independent of $(x,t)\in U\times
[0,T)$. Then for any open subset $V\subset U$ with compact closure
in $U$, there is an extension of the Ricci flow $(V,g(t)|_V)$ past
time $T$.
\end{corollary}

\begin{proof}
\sketch
\uses{thm:shi-derivative-estimates,prop:smooth-patching,def:geometric-limit}
 Take a sequence $t_n\rightarrow T$ and consider the
sequence of Riemannian manifolds $(V,g(t_n))$. By Shi's theorem and the fact
that $V$ has compact closure in $U$, the restriction of this sequence of
metrics to $V$ has uniformly bounded curvature derivatives. Hence, this
sequence has a convergent subsequence with limit $(V,g_\infty)$, where the
convergence is uniform in the $C^\infty$-topology. Now by Hamilton's result
\cite{Hamiltonlimits} it follows that, passing to a further subsequence, the
flows $(V,g(T+t-t_n),(p,0))$ converge to a flow $(V,g_\infty(t),(p,0))$ defined
on $(0,T]$. Clearly, for any $0<t<T$ we have $g_\infty(t)=g(t)$. That is to
say, we have extended the original Ricci flow smoothly to time $T$. Once we
have done this, we extend it to a Ricci flow on $[T,T_1)$ for some $T_1> T$
using the local existence results. The extension to $[T,T_1)$ fits together
smoothly with the flow on $[0,T]$ by Proposition~\ref{prop:smooth-patching}.
\end{proof}




\section{Gromov-Hausdorff convergence}

Let $Z$ be a metric space. Define the
Hausdorff distance between subsets of $Z$ as follows: $d^Z_H(X,Y)$ is the
infimum of all $\delta\ge 0$ such that $X$ is contained in the
$\delta$-neighborhood of $Y$ and $Y$ is contained in the $\delta$-neighborhood
of $X$. For metric spaces $X$ and $Y$ we define the Gromov-Hausdorff distance
between them, denoted $D_{GH}(X,Y)$, to be the infimum over all metric spaces
$Z$ and isometric embeddings $f\colon X\to Z$ and $g\colon Y\to Z$ of the
Hausdorff distance between $f(X)$ and $g(Y)$. For pointed metric spaces $(X,x)$
and $(Y,y)$ of finite diameter, we define the Gromov-Hausdorff
distance between them, denoted
$D_{GH}((X,x),(Y,y))$, to be the infimum of $D_H^Z(f(X),g(Y))$ over all triples
$((Z,z),f,g)$ where $(Z,z)$ is a pointed metric space and $f\colon (X,x)\to
(Z,z)$ and $g\colon (Y,y)\to (Z,z)$ are base-point preserving isometries.


\begin{definition}[$\delta$-net]\label{def:delta-net}
\lean{MorganTianLib.IsDeltaNet}
\leanok
\label{epsnet}
A $\delta$-net in $(X,x)$ is a subset $L$ of $X$ containing $x$
whose $\delta$-neighborhood covers $X$ and for which there is some
$\delta'>0$ with $d(\ell_1,\ell_2)\ge \delta'$ for all $\ell_1\not=
\ell_2$ in $L$.
\end{definition}
 Clearly, the Gromov-Hausdorff distance from a based
metric space $(X,x)$ to a $\delta$-net $(L,x)$ contained in it is at
most $\delta$. Furthermore, for every $\delta>0$ the based space
$(X,x)$ has a $\delta$-net: Consider subsets $L\subset X$ containing
$x$ with the property that the $\delta/2$-balls centered at the
points of $L$ are disjoint. Any maximal such subset (with respect to
the inclusion relation) is a $\delta$-net in $X$.

\begin{lemma}[Triangle inequality for Gromov--Hausdorff distance]\label{lem:gromov-hausdorff-triangle-inequality}
\uses{def:delta-net}
\notready

The Gromov-Hausdorff distance satisfies the triangle inequality.
\end{lemma}

\begin{proof}
\sketch
\uses{def:delta-net}
Suppose that $D_{GH}((X,x),(Y,y))=a$ and $D_{GH}((Y,y),(Z,z))=b$.
Fix any $\delta>0$. Then there is a metric $d_1$ on $X\vee Y$ such
that $d_1$ extends the metrics on $X,Y$ and the
$(a+\delta)$-neighborhood of $X$ is all of $X\vee Y$ as is the
$(a+\delta)$-neighborhood of $Y$. Similarly, there is a metric $d_2$
on $Y\vee Z$ with the analogous properties (with $b$ replacing $a$).
Take a $\delta$-net $(L,y)\subset (Y,y)$, and define
$$d(x',z')=\inf_{\ell\in L}d(x',\ell)+d(\ell,z').$$
We claim that $d(x',z')>0$ unless $x'=z'$ is the common base point. The reason
is that if $\inf_{\ell\in L}d(x',\ell)=0$, then by the triangle
inequality, any sequence of $\ell_n\in L$ with $d(x',\ell_n)$ converging to
zero is a Cauchy sequence, and hence is eventually constant. This means that
for all $n$ sufficiently large, $x'=\ell_n\in L\cap X$ and hence $x'$ is the
common base point. Similarly for $z'$.

A straightforward computation shows that the function $d$ above,
together with the given metrics on $X$ and $Z$, define a metric on
$X\vee Z$ with the property that the $(a+b+3\delta)$-neighborhood of
$X$ is all of $X\vee Z$ and likewise for $Z$. Since we can do this
for any $\delta>0$, we conclude that $D_{GH}((X,x),(Z,z))\le a+b$.
\end{proof}


\begin{definition}[Gromov--Hausdorff convergence (bounded diameter)]\label{def:gromov-hausdorff-convergence}
\uses{def:delta-net}
\notready
\label{GHconverge}

We say that a sequence of  based metric spaces $(X_k,x_k)$ of
uniformly bounded diameter {\sl converges in the Gromov-Hausdorff
sense} to a based metric space $(Y,y)$ of finite
diameter if
$$\lim_{k\rightarrow\infty}D_{GH}((X_k,x_k),(Y,y))=0.$$
Thus, a based metric space $(X,x)$ of bounded diameter is the limit
of a sequence of $\delta_n$-nets $L_n\subset X$ provided that
$\delta_n\rightarrow 0$ as $n\rightarrow\infty$.
\end{definition}

\begin{example}[Shrinking manifolds converge to a point]\label{ex:gromov-hausdorff-limit-point}
\uses{def:gromov-hausdorff-convergence}
\notready

A sequence of compact $n$-manifolds of diameter tending to zero has
a point as  Gromov-Hausdorff limit.
\end{example}


\begin{definition}[Realization sequence]\label{def:realization-sequence}
\uses{def:gromov-hausdorff-convergence}
\notready
\label{seqconv}

Suppose that $\{(X_k,x_k)\}_k$ converges in the Gromov-Hausdorff sense to
$(Y,y)$. Then a {\sl realization sequence} for this convergence is a sequence
of triples $((Z_k,z_k),f_k,g_k)$ where, for each $k$, the pair $(Z_k,z_k)$ is a
based metric space,
$$f_k\colon (X_k,x_k)\to
(Z_k,z_k) \ \ \ \ \text{and} \ \ \ g_k\colon (Y,y)\to (Z_k,z_k)$$ are
isometric embeddings  and $D_{GH}(f_k(X_k),g_k(Y))\rightarrow 0$ as
$k\rightarrow\infty$. Given a realization sequence for the
convergence, we say that a sequence $\ell_k\in X_k$ converges to
$\ell\in Y$ (relative to the given realization sequence) if
$d(f_k(\ell_k),g_k(\ell))\rightarrow 0$ as $i\rightarrow\infty$.
\end{definition}




\begin{lemma}[Net characterization of Gromov--Hausdorff convergence]\label{lem:gromov-hausdorff-net-characterization}
\uses{def:gromov-hausdorff-convergence, def:delta-net}
\notready

The characterization is Proposition 3.5 of \cite{Gromov}.
Let $(X_k,x_k)$ be a sequence of metric spaces whose diameters are uniformly
bounded. Then the $(X_k,x_k)$ converge in the Gromov-Hausdorff
sense to $(X,x)$ if and only if the
following holds for every $\delta>0$. For every $\delta$-net $L\subset X$, for
every $\eta>0$, and for every $k$ sufficiently large, there is a
$(\delta+\eta)$-net $L_k\subset X_k$ and a bijection $L_k\to L$ sending $x_k$
to $x$ so that the push forward of the metric on $L_k$ induced from that of
$X_k$ is $(1+\eta)$-bi-Lipschitz equivalent to the metric on $L$ induced from
$X$.
\end{lemma}


\begin{lemma}[Uniqueness of compact Gromov--Hausdorff limits]\label{lem:gromov-hausdorff-limit-uniqueness-compact}
\uses{def:gromov-hausdorff-convergence}
\notready

Let $(X_k,x_k)$ be a sequence of  based metric spaces whose diameters are
uniformly bounded. Suppose that $(Y,y)$ and $(Y',y')$ are limits in the
Gromov-Hausdorff sense of this sequence and each of $Y$ and $Y'$ are compact.
Then $(Y,y)$ is isometric to $(Y',y')$.
\end{lemma}


\begin{proof}
\sketch
\uses{def:gromov-hausdorff-convergence,lem:gromov-hausdorff-triangle-inequality}
By the triangle inequality for Gromov-Hausdorff distance, it follows
from the hypothesis of the lemma that $D_{GH}((Y,y),(Y',y'))=0$. Fix
$\delta>0$. Since $D_{GH}((Y,y),(Y',y'))=0$, for any $n>0$ and
finite $1/n$-net $L_n\subset Y$ containing $y$ there is an embedding
$\varphi_n\colon L_n\to Y'$ sending $y$ to $y'$ such that the image
is a $2/n$-net in $Y'$ and such that the map from $L_n$ to its image
is a $(1+\delta)$-bi-Lipschitz homeomorphism. Clearly, we can
suppose that in addition the $L_n$ are nested: $L_n\subset
L_{n+1}\subset \cdots$. Since $Y'$ is compact and $L_n$ is finite,
and we can pass to a subsequence so that $\lim_{k\rightarrow\infty}\varphi_k|_{L_n}$ converges to a map
$\psi_n\colon L_n\to Y'$ which is a $(1+\delta)$-bi-Lipschitz map
onto its image which is a $2/n$ net in $Y'$. By a standard
diagonalization argument, we can arrange that
$\psi_{n+1}|_{L_{n}}=\psi_{n}$ for all $n$. The $\{\psi_n\}$ then
define an embedding $\cup_n L_n\to Y'$ that is a
$(1+\delta)$-bi-Lipschitz map onto its image which is a dense subset
of $Y'$. Clearly, using the compactness of $Y'$ this map extends to
a $(1+\delta)$-bi-Lipschitz embedding $\psi_\delta\colon (Y,y)\to
(Y',y')$ onto a dense  subset of $Y'$. Since $Y$ is also compact,
this image is in fact all of $Y'$. That is to say, $\psi_\delta$ is
a $(1+\delta)$-bi-Lipschitz homeomorphism $(Y,y)\to (Y',y')$. Now
perform this construction for  a sequence of $\delta_n\rightarrow 0$
and $(1+\delta_n)$-bi-Lipschitz homeomorphisms
$\psi_{\delta_n}\colon (Y,y)\to (Y',y')$. These form an
equicontinuous family so that by passing to a subsequence we can
extract a limit $\psi\colon (Y,y)\to (Y',y')$. Clearly, this limit
is an isometry.
\end{proof}

\begin{definition}[Gromov--Hausdorff convergence (unbounded diameter)]\label{def:gromov-hausdorff-convergence-unbounded}
\uses{def:gromov-hausdorff-convergence}
\notready

For  based metric spaces $(X_k,x_k)$ (not necessarily of finite
diameter) to {\sl converge in the Gromov-Hausdorff sense to a based
metric space} $(Y,y)$ means
that for each $r>0$ there is a sequence $\delta_k\rightarrow 0$ such
that the sequence of balls $B(x_k,r+\delta_k)$ in $(X_k,x_k)$
converges in the Gromov-Hausdorff sense to the ball $B(y,r)$ in $Y$.
\end{definition}

\begin{lemma}[Uniqueness of locally compact Gromov--Hausdorff limits]\label{lem:gromov-hausdorff-limit-uniqueness-locally-compact}
\uses{def:gromov-hausdorff-convergence-unbounded, lem:gromov-hausdorff-limit-uniqueness-compact}
\notready

Let $(X_k,x_k)$ be a sequence of locally compact metric spaces. Suppose that
$(Y,y)$ and $(Y',y')$ are complete, locally compact, based metric spaces that
are limits of the sequence in the Gromov-Hausdorff sense. Then there is an
isometry $(Y,y)\to (Y',y')$.
\end{lemma}

\begin{proof}
\sketch
\uses{lem:gromov-hausdorff-limit-uniqueness-compact}
We show that for each $r<\infty$ there is an isometry between the
closed balls $\overline{B(y,r)}$ and $\overline{B(y',r)}$. By the
local compactness and completeness,  these closed balls are compact.
Each is the limit in the Gromov-Hausdorff sense of a sequence
$B(x_k,r+\delta_k)$ for some $\delta_k\rightarrow 0$ as
$k\rightarrow\infty$. Thus, invoking the previous lemma we see that
these closed balls are isometric. We take a sequence
$r_n\rightarrow\infty$ and isometries $\varphi_n\colon
(B(y,r_n),y)\to (B(y',r_n),y')$. By a standard diagonalization
argument, we pass to a subsequence such that for each $r<\infty$ the
sequence $\varphi_n|_{B(y,r)}$ of isometry converges to an isometry
$\varphi_r\colon B(y,r)\to B(y',r)$. These then fit together to
define a global isometry $\varphi\colon (Y,y)\to (Y',y')$.
\end{proof}

If follows from this that if a sequence of points $\ell_k\in X_k$ converges to
$\ell\in Y$ under one realization sequence for the convergence and to $\ell'\in
Y$ under another, then there is an isometry of $(Y,y)$ to itself carrying
$\ell$ to $\ell'$.


\begin{example}[Geometric convergence implies Gromov--Hausdorff convergence]\label{ex:geometric-implies-gromov-hausdorff-convergence}
\uses{def:geometric-limit, def:gromov-hausdorff-convergence-unbounded}
\notready

Let $(M_n,g_n,x_n)$ be a sequence of based Riemannian manifolds
converging geometrically to $(M_\infty,g_\infty,x_\infty)$. Then the
sequence also converges in the Gromov-Hausdorff sense to the same
limit.
\end{example}

\subsection{Precompactness}

\begin{definition}[Length space]\label{def:length-space}
\lean{MorganTianLib.LengthSpace}
\leanok
A {\sl length space} is a connected metric space
$(X,d)$ such that for any two points $x,y$ there is a rectifiable
arc $\gamma$ with endpoints $x$ and $y$ and with the length of
$\gamma$ equal to $d(x,y)$.
\end{definition}
For any based metric space $(X,x)$ and constants $\delta>0$ and
$R<\infty$ let $N(\delta,R,X)$ be the maximal number of disjoint
$\delta$-balls in $X$ that can be contained in $B(x,R)$.

\begin{theorem}[Gromov's precompactness theorem]\label{thm:gromov-hausdorff-precompactness}
\uses{def:length-space, def:gromov-hausdorff-convergence-unbounded}
\notready
\label{GHcompact}

This is Gromov's precompactness theorem \cite{Gromov}.
Suppose that $(X_k,x_k)$ is a sequence of based length spaces. Then
there is a based length space $(X,x)$ that is the limit in the
Gromov-Hausdorff sense of a subsequence of the $(X_k,x_k)$ if for
every $\delta>0$ and $R<\infty$ there is an $N<\infty$ such that
$N(\delta,R,X_k)\le N$ for all $k$. On the other hand, if the
sequence $(X_k,x_k)$ has a Gromov-Hausdorff limit, then for every
$\delta>0$ and $R<\infty$ the $N(\delta,R,X_k)$ are bounded
independent of $k$.
\end{theorem}

\subsection{The Tits cone}



Let $(M,g)$ be a complete, non-compact Riemannian manifold of
non-negative sectional curvature. Fix a point $p\in M$, and let
$\gamma$ and $\mu$ be minimal geodesic rays emanating from $p$. For
each $r>0$ let $\gamma(r)$ and $\mu(r)$ be the points along these
geodesic rays at distance $r$ from $p$. Then by  Part 1 of
Theorem~\ref{thm:length-comparison} we see that
$$\ell(\gamma,\mu,r)=\frac{d(\gamma(r),\mu(r))}{r}$$
is a non-increasing function of $r$. Hence, there is a limit
$\ell(\gamma,\mu)\ge 0$ of $\ell(\gamma,\mu,r)$ as
$r\rightarrow\infty$. We define the {\sl angle at infinity} between
$\gamma$ and $\mu$, $0\le\theta_\infty(\gamma,\mu)\le \pi$, to be
the angle at $b$ of the Euclidean triangle $a,b,c$ with side lengths
$|ab|=|bc|=1$ and $|bc|=\ell(\gamma,\mu)$. If $\nu$ is a third geodesic ray
emanating from $p$, then clearly,
$\theta_\infty(\gamma,\mu)+\theta_\infty(\mu,\nu)\ge
\theta_\infty(\gamma,\nu)$.



\begin{definition}[Space of directions at infinity]\label{def:space-of-directions-infinity}
\uses{def:length-space, thm:length-comparison}
\notready

Now we define a metric space whose underlying space is the quotient
space of the equivalence classes of minimal geodesic rays emanating
from $p$, with two rays equivalent if and only if the angle at
infinity between them is zero. The pseudo-distance function
$\theta_\infty$ descends to a metric on this space. This space is a
 length space \cite{BGP}. Notice that the distance between
any two points in this metric space is at most $\pi$. We denote this
space by $S_\infty(M,p)$.
\end{definition}

\begin{claim}[Space of directions at infinity is compact]\label{claim:space-of-directions-compact}
\uses{def:space-of-directions-infinity}
\notready

$S_\infty(M,p)$ is a compact space.
\end{claim}

\begin{proof}
\sketch
Let $\{[\gamma_n]\}_n$ be a sequence of points in $S_\infty(M,p)$. We show that
there is a subsequence with a limit point. By passing to a subsequence we can
arrange that the unit tangent vectors to the $\gamma_n$ at $p$ converge to a
unit tangent vector $\tau$, say. Fix $d<\infty$, and let $x_n$ be the point of
$\gamma_n$ at distance $d$ from $p$. Then by passing to a subsequence we can
arrange that the $x_n$ converge to a point $x$. The minimizing geodesic
segments $[p,x_n]$ on $\gamma_n$ then converge to a minimizing geodesic segment
connecting $p$ to $x$. Performing this construction for a sequence of $d$
tending to infinity and then taking a diagonal subsequence produces a
minimizing geodesic ray $\gamma$ from $p$ whose class is the limit of a
subsequence of the $\{[\gamma_n]\}$
\end{proof}

 We define the {\sl Tits cone} of $M$ at $p$, denoted ${\mathcal T}(M,p)$,
 to be the cone over
$S_\infty(M,p)$, i.e.,  the quotient of the space $S_\infty(M,p)\times
[0,\infty)$ where all points $(x,0)$ are identified together (to become the
cone point). The cone metric on this space is given as follows: Let $(x_1,a_1)$
and $(x_2,a_2)$ be points of $S_\infty(M,p)\times [0,\infty)$. Then the
distance between their images in the cone is determined by
$$d^2([x_1,a_1],[x_2,a_2])=a_1^2+a_2^2-2a_1a_2\cos(\theta_\infty(x_1,x_2)).$$
 It is an easy exercise to show that the Tits cone of $M$ at $p$ is in fact
independent of the choice of $p$. From the previous claim, it follows that the
Tits cone of $M$ is locally compact and complete.

\begin{proposition}[Tits cone as a Gromov--Hausdorff limit]\label{prop:tits-cone-gromov-hausdorff-limit}
\uses{def:space-of-directions-infinity, claim:space-of-directions-compact, def:gromov-hausdorff-convergence-unbounded}
\notready
\label{conelimit}

Let $(M,g)$ be a non-negatively curved, complete, non-compact
Riemannian manifold of dimension $k$. Fix a point $p\in M$ and let
$\{x_n\}_{n=1}^\infty$ be a sequence tending to infinity in $M$. Let
$\lambda_n=d^2(p,x_n)$ and consider the sequence of based Riemannian
manifolds $(M,g_n,p)$, where $g_n=\lambda_n^{-1}g$. Then there is a
subsequence converging in the Gromov-Hausdorff sense. Any
Gromov-Hausdorff limit
 $(X,g_\infty,x_\infty)$  of a
subsequence $(X,g_\infty)$ is isometric to the Tits cone ${\mathcal
T}(M,p)$ with base point the cone point.
\end{proposition}

\begin{proof}
\sketch
\uses{thm:length-comparison}
Let $c$ be the cone point of ${\mathcal T}(M,p)$, and denote by $d$ the
distance function on ${\mathcal T}(M,p)$. Consider the ball $B(c,R)\subset
{\mathcal T}(M,p)$. Since $S_\infty(M,p)$ is the metric completion of the
quotient space of minimal geodesic rays emanating from $p$, for any $\delta>0$
there is a $\delta$-net $L\subset B(c,R)$ consisting of the cone point together
with points of the form $([\gamma],t)$ where $\gamma$ is a minimal geodesic ray
emanating from $p$ and $t>0$. We define a map from $\psi_n\colon L \to (M,g_n)$
by sending the cone point to $p$ and sending $([\gamma],t)$ to the point at
$g_n$-distance $t$ from $p$ along $\gamma$. Clearly, $\psi_n(L)$ is contained
in $B_{g_n}(p,R)$. From the second item of Theorem~\ref{thm:length-comparison} and the
monotonicity of angles it follows that the map $\psi_n\colon L\to (M,g_n)$ is a
distance non-decreasing map; i.e., $\psi_n^*(g_n|\psi_n(L))\ge d|_L$. On the
other hand, by the monotonicity, $\psi_{n+1}^*(g_{n+1}|\psi_{n+1}(L))\le
\psi_n^*(g_n|\psi_n(L))$ and this non-increasing sequence of metrics converges
to $d|_L$. This proves that for any $\delta>0$ for all $n$ sufficiently large,
the embedding $\psi_n$ is a $(1+\delta)$-bi-Lipschitz homeomorphism.

 It remains to show that for any $\eta>0$ the images $\psi_n(L)$ are eventually
$\delta+\eta$-nets in $B_{g_n}(p,R)$. Suppose not. Then after
passing to a subsequence, for each $n$ we have a point $x_n\in
B_{g_n}(p,R)$ whose distance from $\psi_n(L)$ is at least
$\delta+\eta$. In particular, $d_{g_n}(x_n,p)\ge \delta$. Consider a
sequence of minimal geodesic rays $\mu_n$ connecting $p$ to the
$x_n$. Since the $g$-length of $\mu_n$ is at least $n\delta$,  by
passing to a further subsequence, we can arrange that the $\mu_n$
converge to a minimal geodesic ray $\gamma$ emanating from $p$. By
passing to a further subsequence if necessary, we arrange that
$d_{g_n}(x_n,p)$ converges to $r>0$. Now consider the points $\tilde
x_n$ on $\gamma$ at $g$-distance $\sqrt{\lambda_n}r$ from $p$.
Clearly, from the second item of Theorem~\ref{thm:length-comparison} and the
fact that the angle at $p$ between the $\mu_n$ and $\mu$ tends to
zero as $n\rightarrow\infty$ we have $d_{g_n}(x_n,\tilde
x_n)\rightarrow 0$ as $n\rightarrow\infty$. Hence, it suffices to
show that for all $n$ sufficiently large, $\tilde x_n$ is within
$\delta$ of $\psi_n(L)$ to obtain a contradiction. Consider the
point $z=([\mu],r)\in {\mathcal T}(M,p)$. There is a point
$\ell=([\gamma],t')\in L$ within distance $\delta$ of $z$ in the
metric $d$. Let $\tilde y_n\in M$ be the point in $M$ at
$g$-distance $\sqrt{\lambda_n}t'$ along $\gamma$. Of course, $\tilde
y_n=\psi_n(\ell)$. Then $d_{g_n}(\tilde x_n,\tilde y_n)\rightarrow
d(\ell,z)<\delta$. Hence, for all $n$ sufficiently large,
$d_{g_n}(\tilde x_n,\tilde y_n)<\delta$. This proves that for all
$n$ sufficiently large $\tilde x_n$ is within $\delta$ of
$\psi_n(L)$ and hence for all $n$ sufficiently large $x_n$ is within
$\delta+\eta$ of $\psi_n(L)$.

We have established that for every $\delta,\eta>0$ and every $R<\infty$ there
is a finite $\delta$-net $L$ in $({\mathcal T}(M,p),c)$ and for all $n$
sufficiently large an $(1+\delta)$-bi-Lipschitz embedding $\psi_n$ of $L$ into
$(M,g_n,p)$ with image a $\delta+\eta$-net for $(M,g_n,p)$. This proves that
the sequence $(M,g_n,p)$ converges in the Gromov-Hausdorff sense to ${\mathcal
T}(M,p),c))$.
\end{proof}


\section{Blow-up limits}

Here we introduce a type of geometric limit. These were  originally
introduced and studied by Hamilton in \cite{Hamiltonsurvey}, where,
among other things, he showed that  $3$-dimensional blow-up limits
have non-negative sectional curvature. We shall use repeatedly
blow-up limits and the positive curvature result in the arguments in
the later sections.

\begin{definition}[Blow-up limit]\label{def:blow-up-limit}
\uses{def:geometric-limit, def:generalized-ricci-flow}
\notready

Let $({\mathcal M}_k,G_k,x_k)$ be a sequence of based
generalized Ricci flows. We suppose that ${\bf t}(x_k)=0$ for all
$n$. We set $Q_k$ equal to $R(x_k)$. We denote by $(Q_k{\mathcal
M}_k,Q_kG_k,x_k)$ the family of generalized flows that have been
rescaled so that $R_{Q_kG_k}(x_k)=1$.
 Suppose that $\lim_{k\rightarrow\infty}Q_k=\infty$ and that after passing to a
subsequence there is a geometric limit of the sequence
$(Q_k{\mathcal M}_k,Q_kG_k,x_k)$ which is a Ricci flow defined for
$-T<t\le 0$. Then we call this limit a {\em blow-up
limit} of the original based sequence. In
the same fashion, if there is a geometric limit for a subsequence of
the zero time-slices of the $(Q_k{\mathcal M}_k,Q_kG_k,x_k)$, then
we call this limit
 the  blow-up limit of the $0$ time-slices.
\end{definition}

The significance of the condition that the generalized Ricci flows
have curvature pinched toward positive is that, as Hamilton
originally established in \cite{Hamiltonsurvey}, the latter
condition implies that any blow-up limit has non-negative curvature.

\begin{theorem}[Blow-up limits have non-negative curvature]\label{thm:blow-up-limit-nonnegative-curvature}
\uses{def:blow-up-limit}
\notready
\label{blowupposcurv}

Let $({\mathcal M}_k,G_k,x_k)$ be a sequence of generalized $3$-dimensional
Ricci flows, each of which has time interval of definition contained in
$[0,\infty)$ and each of which has curvature pinched toward positive. Suppose
that $Q_k=R(x_k)$ tends to infinity as $k$ tends to infinity. Let $t_k={\bf
t}(x_k)$ and let $({\mathcal M}_k',G_k',x_k)$ be the result of shifting time by
$-t_k$ so that ${\bf t}'(x_k)=0$. Then any blow-up limit of the sequence
$({\mathcal M}_k,G_k',x_k)$ has non-negative Riemann curvature. Similarly, any
blow-up limit of the zero time-slices of this sequence has non-negative
curvature.
\end{theorem}

\begin{proof}
\sketch
\uses{def:blow-up-limit}
Let us consider the case of the geometric limit of the zero time-slice first.
Let $( M_\infty,g_\infty(0),x_\infty)$ be a blow-up limit of the zero
time-slices in the sequence. Let $V_k\subset M_\infty$ and $\varphi_k\colon
V_k\to (M_k)_0$ be as in the definition of the geometric limit. Let $y\in
{\mathcal M}_\infty$ be a point and let $\lambda(y)\ge \mu(y)\ge \nu(y)$ be the
eigenvalues of the Riemann curvature operator for $g_\infty$ at $y$. Let
$\{y_k\}$ be a sequence in $Q_k{\mathcal M}_k'$ converging to $y$, in the sense
that $y_k=\varphi_k(y)$ for all $k$ sufficiently large. Then
\begin{eqnarray*}
\lambda(y) & = & \lim_{n\rightarrow\infty}Q_k^{-1}\lambda(y_k) \\
\mu(y) & = & \lim_{n\rightarrow\infty}Q_k^{-1}\mu(y_k) \\
\nu(y) & = & \lim_{n\rightarrow\infty}Q_k^{-1}\nu(y_k)
\end{eqnarray*}
Since by Equation~(\ref{Rlower}) we have $R(y_k)\ge -6$ for all $k$ and since
by hypothesis $Q_k$ tends to infinity as $n$ does, it follows that $R(y)\ge 0$.
Thus if $\lambda(y)=0$, then $\Rm(y)=0$ and the result is established at
$y$. Hence, we may assume that $\lambda(y)>0$, which means that $\lambda(y_k)$
tends to infinity as $k$ does. If $\nu(y_k)$ remains bounded below as $k$ tends
to infinity, then $Q_k^{-1}\nu(y_k)$ converges to a limit which is $\ge 0$, and
consequently $Q_k^{-1}\mu(y_k)\ge Q_k^{-1}\nu(y_k)$ has a non-negative limit.
Thus, in this case the Riemann curvature of $g_\infty$ at $y$ is non-negative.
On the other hand, if $\nu(y_k)$ goes to $-\infty$ as $k$ does, then according
to Equation~(\ref{RXineq}) the ratio of $X(y_k)/R(y_k)$ goes to zero. Since
$Q_k^{-1}R(y_k)$ converges to the finite limit $R(y)$, the product
$Q_k^{-1}X(y_k)$ converges to zero as $k$ goes to infinity. This means that
$\nu(y)=0$ and consequently that $\mu(y)\ge 0$. Thus, once again we have
non-negative curvature for $g_\infty$ at $y$.

The argument in the case of a geometric limit flow is identical.
\end{proof}

\begin{corollary}[Blow-up limits of pinched flows have non-negative curvature]\label{cor:blow-up-limit-nonnegative-curvature-pinched}
\uses{thm:blow-up-limit-nonnegative-curvature}
\notready

Suppose that $(M_k,g_k(t))$ is a sequence of Ricci flows each of which has time
interval of definition contained in $[0,\infty)$ with each $M_k$ being a
compact $3$-manifold. Suppose further that, for each $k$, we have $|\Rm(p_k,0)|\le 1$ for all $p_k\in M_k$. Then any blow-up limit of this sequence
of Ricci flows has non-negative curvature.
\end{corollary}


\begin{proof}
\sketch
\uses{thm:blow-up-limit-nonnegative-curvature,thm:pinching-toward-positive-curvature}
According to Theorem~\ref{pinch} the hypotheses imply that for every $k$ the
Ricci flow $(M_k,g_k(t))$ has curvature pinched toward positive. From this, the
corollary follows immediately from the previous theorem.
\end{proof}

\section{Splitting limits at infinity}

In our later arguments we shall need a splitting result at infinity
in the non-negative curvature case. Assuming that a geometric limit
exists, the splitting result is quite elementary. For this reason we
present it here, though it will not be used until
Chapter~\ref{kappasect}.



The main result of this section gives a condition under which a geometric limit
automatically splits off a line.

\begin{theorem}[Splitting at infinity]\label{thm:splitting-at-infinity}
\uses{def:geometric-limit}
\notready
\label{topsplit}

Let $(M,g)$ be a complete, connected manifold of non-negative
sectional curvature. Let $\{x_n\}$ be a sequence of points going off
to infinity, and suppose that we can find scaling factors
$\lambda_n>0$ such that the based Riemannian manifolds
$(M,\lambda_ng,x_n)$ have a geometric limit
$(M_\infty,g_\infty,x_\infty)$. Suppose that there is a point $p\in
M$ such that $\lambda_nd^2(p,x_n)\rightarrow \infty$ as
$n\rightarrow\infty$.  Then, after passing to a subsequence,
minimizing geodesic arcs $\gamma_n$ from $x_n$ to $p$ converge to a
minimizing geodesic ray in $M_\infty$. This minimizing geodesic ray
is part of a minimizing geodesic line $\ell$ in $M_\infty$. In
particular, there is a Riemannian product decomposition
$M_\infty=N\times \Ar$ with the property that $\ell$ is $\{x\}\times
\Ar$ for some $x\in N$.
\end{theorem}


\begin{proof}
\sketch
\uses{thm:length-comparison,lem:minimizing-line-implies-product}
Let $d_n$ be the distance from $p$ to $x_n$. Consider minimizing
geodesic arcs $\gamma_n$ from $p$ to $x_n$. By passing to a
subsequence we can assume that tangent directions at $p$ of these
arcs converge. Hence, for every $0<\delta<1$ there is $N$ such that
for all $n,m\ge N$ the angle between $\gamma_n$ and $\gamma_m$ at
$p$ is less than $\delta$. For any $n$ we can choose $m(n)$ such
that $d_{m(n))}\ge d_n(1+1/\delta)$. Let $\mu_n$ be a minimizing
geodesic from $x_n$ to $x_{m(n)}$.
 Now applying the Toponogov comparison
(first part of Theorem~\ref{thm:length-comparison}) and the usual law of
cosines in Euclidean space, we see that the distance $d$ from $x_n$
to $x_{m(n)}$ satisfies
$$d_{m(n)}-d_n\le d\le \sqrt{d_n^2+d_{m(n)}^2-2d_nd_{m(n)}\cos(\delta)}.$$ Let $\theta_n=\angle_{x_n'}$ of the Euclidean triangle
$\triangle(x_n',p',x_{m(n)}')$ with $|s_{x_n'p'}|=d_n,|s_{x_n'x_{m(n)}'}|=d$
and $|s_{p'x_{m(n)}'}|=d_{m(n)}$. Then for any $\alpha<d_n$ and $\beta<d$ let
$x$ and $y$ be the points on $s_{x_np}$ and on $s_{x_nx_{m(n)}}$ at distances
$\alpha$ and $\beta$ respectively from $x_n$. Given this, according to the
Toponogov comparison result (first part of Theorem~\ref{thm:length-comparison}), we have
$$d(x,y)\ge \sqrt{\alpha^2+\beta^2-2\alpha\beta\cos(\theta_n)}.$$
The angle $\theta_n$ satisfies:
$$d_n^2+d^2-2d_nd\cos(\theta_n)=d_{m(n)}^2.$$
Thus, \begin{eqnarray*} \cos(\theta_n) & = &
\frac{d_n^2+d^2-d_{m(n)}^2}{2d_nd} \\
& \le & \frac{2d_n^2-2d_nd_{m(n)}\cos(\delta)} {2d_nd} \\
& = & \frac{d_n}{d}-\frac{d_{m(n)}}{d}\cos(\delta)\\
& \le & \delta-(1-\delta)\cos(\delta).
\end{eqnarray*}
Since $\delta\rightarrow 0$ as $n\rightarrow\infty$, it follows that
given any $\delta>0$, for all  $n$ sufficiently large, $1+\cos(\theta_n)<\delta$.

We are assuming that the based Riemannian manifolds
$\{(M,\lambda_ng,x_n)\}_{n=1}^\infty$ converge to a geometric limit
$(M_\infty,g_\infty,x_\infty)$. Also, by assumption,
$d_{\lambda_ng_n}(p,x_n)\rightarrow \infty$ as $n\rightarrow
\infty$, so that the lengths of the $\gamma_n$ tend to infinity in
the metrics $\lambda_ng_n$. This also means that the lengths of
$\mu_n$, measured in the metrics $\lambda_ng_n$, tend to infinity.
Thus, by passing to a subsequence we can assume that each of these
families, $\{\gamma_n\}$ and $\{\mu_n\}$, of minimizing geodesic
arcs converges to a minimizing geodesic arc, which we denote $\tilde
\gamma$ and $\tilde \mu$, respectively, in $M_\infty$ emanating from
$x_\infty$. The above computation shows that the angle between these
arcs is $\pi$ and hence that their union is a geodesic, say $\ell$.
The same computation shows that  $\ell $ is minimizing.

The existence of the minimizing geodesic line $\ell$ together with
the fact that the sectional curvatures of the limit are $\ge 0$
implies by Lemma~\ref{lem:minimizing-line-implies-product} that the limit manifold is a Riemannian
product $N\times \Ar$ in such a way that $\ell$ is of the form
$\{x\}\times \Ar$ for some $x\in N$.
\end{proof}
